\chapter{Descripción de sistemas}
\label{cap:sistemas}

\opening{Comprender un helicóptero es comprender cómo cada sistema depende de los demás.}

\section{Generalidades}
\fuente{POH, sección 7, con las revisiones del 22 FEB 2024}

El R22 es un helicóptero biplaza, de rotor principal único y motor único, construido principalmente en metal y equipado con tren de aterrizaje de patines.

La estructura primaria del fuselaje es de tubos de acero soldados y chapa de aluminio remachada. El cono de cola es una estructura monocasco en la que los revestimientos de aluminio soportan las cargas primarias. La fibra de vidrio y los termoplásticos se usan en la estructura secundaria de la cabina, los carenados de enfriamiento del motor y diversos ductos y carenados. Las puertas de la cabina también son de fibra de vidrio y termoplásticos.

Una puerta de capó (\emph{cowl door}) en el costado derecho da acceso a la caja de transmisión principal y al sistema de transmisión. Paneles y capós removibles dan acceso adicional a los controles y otros componentes para el mantenimiento. Hay mamparos cortafuego de acero inoxidable delante y encima del motor. Ambas puertas de la cabina son removibles (vea el capítulo~\ref{cap:manejo}).

\section{Sistemas rotores}

El \textbf{rotor principal} tiene dos palas totalmente metálicas montadas en el cubo mediante bisagras de conicidad (\emph{coning hinges}). El cubo se monta en el eje mediante una bisagra de balanceo (\emph{teeter hinge}). Ambas bisagras usan rodamientos autolubricados. Los topes de caída (\emph{droop stops}) de las palas proporcionan una restricción por fricción en la bisagra de balanceo que normalmente impide que el rotor balancee durante la detención o la puesta en marcha.

Los rodamientos de cambio de paso de cada pala están encerrados en una carcasa en la raíz de la pala, llena de aceite y sellada con una funda elastomérica. Cada pala tiene un larguero grueso de acero inoxidable en el borde de ataque, resistente a la corrosión y a la erosión. Los revestimientos están pegados al larguero aproximadamente una pulgada detrás del borde de ataque. \textbf{Las palas deben repintarse si la pintura se erosiona hasta el metal desnudo en la línea de unión revestimiento--larguero}: la unión puede dañarse si la línea queda expuesta.

El \textbf{rotor de cola} tiene dos palas totalmente metálicas y un cubo balancín con ángulo de conicidad fijo. Los rodamientos de cambio de paso tienen camisas autolubricadas. Los rodamientos de la bisagra de balanceo son elastoméricos o tienen camisas autolubricadas. Las palas tienen revestimientos y herrajes de raíz de aluminio. Mantener la pintura en buen estado reduce la corrosión y la erosión.

\section{Sistema de transmisión}

Una polea para correas en V está apernada directamente al eje de salida del motor. Las correas transmiten la potencia a la polea superior, que lleva en su cubo un embrague de rueda libre. El eje interior del embrague transmite la potencia hacia adelante, al rotor principal, y hacia atrás, al rotor de cola. Hay acoplamientos flexibles en la entrada de la caja principal y en cada extremo del eje largo del rotor de cola.

La caja de transmisión principal contiene un único juego de engranajes cónicos espirales de una etapa, lubricado por salpicadura. Un ducto de enfriamiento bajo la caja está conectado a la parte superior del carenado del motor. La caja principal se apoya en cuatro soportes de goma.

El eje largo del rotor de cola no tiene rodamientos de apoyo, pero sí un rodamiento amortiguador poco cargado. La caja de cola contiene un único juego de engranajes cónicos espirales a 90°, lubricado por salpicadura.

\section{Grupo motor}

Un motor Lycoming de cuatro cilindros horizontalmente opuestos, válvulas en cabeza, enfriado por aire, con carburador y sistema de aceite de cárter húmedo, propulsa el helicóptero. El motor está equipado con motor de arranque, alternador, encendido blindado, dos magnetos, silenciador, enfriador de aceite y filtro de aire de inducción.

Un ventilador tipo jaula de ardilla de accionamiento directo, montado en el eje de salida del motor, suministra aire de enfriamiento a los cilindros y al enfriador de aceite a través de un carenado de fibra de vidrio y aluminio.

El aire de inducción entra por una toma en el costado derecho del fuselaje y pasa por un ducto flexible a la caja de aire del carburador. Una toma de aire caliente suministra aire calentado a la caja. Una válvula deslizante controlada por el piloto deja pasar aire frío o caliente a la caja, a través del filtro de aire y hacia el carburador.

El piloto debe leer y respetar los procedimientos recomendados en el manual del operador de Lycoming para obtener la máxima vida y eficiencia del motor.

\section{Controles de vuelo}

Los dobles comandos son equipo estándar y todos los controles primarios se accionan mediante tubos de empuje--tracción y balancines (\emph{bellcranks}). Los rodamientos del sistema de control son rodamientos de bolas sellados que no requieren lubricación o tienen camisas autolubricadas.

El funcionamiento de los controles es convencional. El cíclico está montado al centro, con las empuñaduras izquierda y derecha montadas en un tubo transversal que pivota sobre el poste central del cíclico. En los helicópteros más recientes, un mecánico puede ajustar el ángulo de la empuñadura del piloto hacia adelante o hacia atrás respecto al tubo transversal para lograr la posición más cómoda de la mano. La posición más adelantada da el mayor espacio de control con cíclico atrás para pilotos de mayor tamaño. Los pilotos siempre deben verificar antes del vuelo que pueden aplicar todo el recorrido de los controles.

El colectivo también es convencional. El acelerador del motor está correlacionado con el colectivo mediante una articulación mecánica: al subir el colectivo se abre el acelerador y al bajarlo se cierra. El colectivo incorpora además un control de acelerador de empuñadura giratoria (vea ``Controles del motor'').

\begin{precaucion}
Sobre 4000~ft, la correlación acelerador--colectivo y el gobernador son menos eficaces. Por lo tanto, los cambios de potencia deben ser lentos y suaves.
\end{precaucion}

\begin{precaucion}
Con potencias altas sobre 4000~ft, el acelerador suele estar completamente abierto y la RPM debe controlarse con el colectivo.
\end{precaucion}

\section{Controles de vuelo removibles}

Los controles del asiento izquierdo pueden ser removidos e instalados por el personal de mantenimiento o por los pilotos:

\begin{pasos}
\item \textbf{Empuñadura del cíclico:} retire el pasador de liberación rápida presionando el botón y tirando, y luego tire hacia afuera de la empuñadura izquierda sosteniendo el poste central del cíclico. Para instalarla, invierta el procedimiento.
\end{pasos}

\begin{nota}
Los helicópteros más recientes tienen un anillo moleteado junto al pasador de liberación rápida que puede apretarse a mano para eliminar el juego. El anillo debe estar suelto (gírelo en sentido antihorario mirando hacia el centro) para retirar el pasador.
\end{nota}

\begin{precaucion}
Después de retirar la empuñadura del cíclico, coloque la tapa plástica en el extremo expuesto del tubo transversal para evitar posibles lesiones al ocupante del asiento izquierdo.
\end{precaucion}

\begin{pasos}
\setcounter{enumi}{1}
\item \textbf{Colectivo:} empuje la funda hacia atrás para dejar a la vista los pasadores de traba. Presione los pasadores y tire del bastón hacia adelante. Para instalarlo, invierta el procedimiento; puede ser necesario girar levemente el bastón para que los pasadores encajen.
\end{pasos}

\begin{precaucion}
Con el colectivo instalado, asegúrese de que ambos pasadores de traba atraviesen completamente los agujeros de cada lado.
\end{precaucion}

\begin{pasos}
\setcounter{enumi}{2}
\item \textbf{Pedales del rotor de cola:} presione el pasador de traba mientras gira el pedal en sentido antihorario y tire hacia arriba. Para instalarlos, invierta el procedimiento.
\end{pasos}

\section{Gobernador de RPM}

El gobernador mantiene la RPM del motor detectando sus variaciones y aplicando correcciones al acelerador mediante un embrague de fricción que el piloto puede vencer con facilidad. El gobernador está activo solo sobre 80\,\% de RPM del motor y se conecta o desconecta con el interruptor de palanca en el extremo del colectivo derecho.

El gobernador está diseñado para ayudar a controlar la RPM en condiciones normales. Puede no evitar sobrevelocidades o bajas de RPM generadas por maniobras agresivas.

\begin{precaucion}
Al operar con altitudes de densidad altas, la respuesta del gobernador puede ser demasiado lenta para evitar una sobrevelocidad durante ráfagas, tirones o al bajar el colectivo.
\end{precaucion}

\section{Compensación y fricción de los controles}

El cíclico y el colectivo incorporan resortes de compensación. El resorte de colectivo arriba equilibra las cargas del rotor y permite al piloto soltar la mano izquierda del colectivo en la mayoría de los regímenes de vuelo. El cíclico longitudinal tiene un resorte (\emph{bungee}) fijo que anula la mayor parte de las fuerzas longitudinales en crucero.

El cíclico lateral tiene un resorte de compensación de encendido/apagado que anula la fuerza hacia la izquierda que aparece a alta velocidad. Se acciona con un pomo de empujar--tirar ubicado justo delante del cíclico. Desde el S/N 550, el ajuste fino de la fuerza se controla con el pomo del costado izquierdo de la consola.

\begin{precaucion}
Si se tira inadvertidamente del control de mezcla en vuelo, el motor se detendrá. Para no tirar del control equivocado, acceda siempre a la compensación lateral pasando la mano por el lado izquierdo del poste central del cíclico.
\end{precaucion}

El cíclico y el colectivo tienen dispositivos de fricción ajustables. La palanca de fricción del colectivo está cerca del extremo trasero del colectivo del piloto: hacia atrás aumenta la fricción y hacia adelante la libera. El pomo de fricción del cíclico está a la izquierda del poste central: girarlo en sentido horario aplica fricción al cíclico longitudinal y lateral.

\begin{precaucion}
La fricción de los controles debe usarse con cuidado en vuelo. Una fricción excesiva puede hacer difícil controlar el helicóptero.
\end{precaucion}

Los pedales accionan mandos de empuje--tracción conectados directamente al control de paso del rotor de cola y no tienen resorte de compensación ni fricción.

\section{Controles del motor}

Cada colectivo tiene un control de acelerador de empuñadura giratoria. Los controles están interconectados y accionan la válvula del acelerador mediante una articulación mecánica. El acelerador también está correlacionado con el colectivo: al subirlo se abre y al bajarlo se cierra.

El gobernador electrónico hace pequeños ajustes del acelerador girando la empuñadura para mantener la RPM dentro de los límites con potencia. Normalmente no es necesario manipular la empuñadura salvo en el arranque, la detención, las autorrotaciones de práctica y las emergencias.

Un resorte de sobrerrecorrido en la articulación del acelerador permite al piloto girarlo más allá del tope de ralentí antes de un aterrizaje en autorrotación con contacto (\emph{run-on}). Esto impide que el acelerador se abra al subir el colectivo.

El ajuste correcto de la articulación puede verificarse durante el prevuelo girando la empuñadura a través del resorte de sobrerrecorrido y sosteniéndola contra el tope duro de ralentí: el brazo del acelerador del carburador apenas debe empezar a moverse cuando el colectivo se sube por completo.

Otros controles del motor son el control de mezcla, delante y a la derecha del poste central del cíclico, y el control de calefacción del carburador, a la derecha y detrás del cíclico. Los R22 con motor O-360 tienen el asistente de calefacción del carburador descrito en el capítulo~\ref{cap:normales}.

\begin{precaucion}
No se recomienda empobrecer la mezcla en vuelo. Si se empobrece en tierra a gran altitud, asegúrese de volver a enriquecerla antes de descender a una altitud menor; de lo contrario, el motor puede detenerse.
\end{precaucion}

\begin{nota}
En algunos helicópteros el control de mezcla está en la cara de la consola. Con este control no se usa el protector de mezcla.
\end{nota}

\section{Actuador del embrague}

Después de arrancar, el motor se acopla al sistema de transmisión del rotor mediante las correas en V, que se tensan elevando la polea superior. Un actuador eléctrico, ubicado entre las poleas, eleva la polea superior cuando el piloto acciona el interruptor del embrague. El actuador detecta la carga de compresión (la tensión de las correas) y se apaga cuando las correas quedan correctamente tensadas. La luz de precaución del embrague se enciende siempre que el circuito del actuador está energizado ---acoplando, desacoplando o retensando--- y permanece encendida hasta que las correas están correctamente tensadas o completamente desacopladas.

La holgura de las correas durante el arranque debe ajustarse de modo que las palas empiecen a girar dentro de cinco segundos de acoplar el embrague. Una holgura excesiva puede hacer que las correas salten fuera de las ranuras de las poleas durante el arranque. A medida que las correas se desgastan en servicio, un mecánico puede tener que reajustarlas periódicamente.

Un fusible ubicado en el panel de interruptores de prueba, o cerca de él, impide que una sobrecarga del motor del actuador dispare el disyuntor. Si el fusible se funde, el motor del actuador se detiene pero la luz del embrague permanece encendida. Un disyuntor abierto quita la alimentación tanto al motor como a la luz: no habrá tensado de correas y la luz no podrá indicar una condición anormal.

\begin{precaucion}
Nunca despegue con la luz de precaución del embrague encendida.
\end{precaucion}

\section{Sistema de combustible}

El sistema de combustible es por gravedad (sin bombas) e incluye un tanque principal, un tanque auxiliar opcional, un control de la válvula de corte ubicado detrás del asiento izquierdo y un filtro decantador (\emph{gascolator}). En los helicópteros más recientes los tanques tienen vejigas flexibles dentro de recintos de aluminio; los anteriores usan tanques todo aluminio. Los respiraderos de los tanques están dentro del carenado del mástil.

Hay válvulas de drenaje tipo émbolo para el gascolator y para el sumidero de cada tanque. El gascolator está en la parte inferior izquierda del cortafuego. Los drenajes del tanque auxiliar y del tanque principal con vejiga están dentro de la puerta del capó, bajo el tanque auxiliar; tubos plásticos unidos a las válvulas permiten drenar el combustible hacia afuera. Las muestras se toman presionando el émbolo. En los tanques principales todo aluminio, el drenaje está en el costado izquierdo del fuselaje y se abre empujando el tubo plástico. En los helicópteros nuevos, un tubo de vidrio guardado dentro de la puerta superior trasera del capó sirve para recoger las muestras. \textbf{Debe tomarse muestra de las tres ubicaciones antes del primer vuelo del día y después de cada carga de combustible}, para verificar que no haya contaminación y que el grado sea el correcto.

Los indicadores de combustible funcionan eléctricamente con transmisores de flotador en los tanques. Cuando marcan E (vacío), los tanques están vacíos salvo una pequeña cantidad de combustible no utilizable. La luz de bajo combustible la acciona un sensor eléctrico independiente en el fondo del tanque principal.

El tanque auxiliar está interconectado con el principal y ubicado algo más arriba, de modo que se vacía primero mientras aún queda combustible en el principal. La válvula de corte controla el flujo de ambos tanques al motor.

\section{Sistema eléctrico}

El sistema estándar es de 14~V de corriente continua, con alternador y batería sellada de plomo--ácido. La batería está en el compartimiento del motor o bajo la consola de instrumentos.

El panel de disyuntores está en la repisa justo delante del asiento izquierdo; cada disyuntor está rotulado con su función y amperaje. No se recomienda reponer disyuntores en vuelo.

El interruptor de batería controla el relé de batería, que desconecta la batería del sistema. Un cable protegido por un fusible cerca de la batería puentea el relé para que ambos tacómetros y el reloj sigan recibiendo energía de la batería con el interruptor de batería apagado.

La unidad de control del alternador protege el sistema de las sobretensiones. El amperímetro indica la corriente hacia la batería (``$-$'' indica descarga). Una luz \textsf{ALT} o una indicación de descarga en el amperímetro en vuelo indica bajo voltaje y una posible falla del alternador. Apague el equipo no esencial y ponga el alternador en OFF y de nuevo en ON después de un segundo para reiniciar la unidad de control. Si la luz sigue encendida o el amperímetro sigue indicando descarga, aterrice tan pronto como sea práctico.

\begin{precaucion}
Continuar el vuelo sin alternador funcionando puede producir la pérdida de alimentación de los tacómetros, una condición de vuelo peligrosa.
\end{precaucion}

\begin{nota}
Salvo en los procedimientos de emergencia, no opere el alternador con la batería apagada. La batería ayuda a proteger el equipo eléctrico de los picos de tensión.
\end{nota}

Los helicópteros más recientes tienen un interruptor maestro de aviónica que controla la alimentación de la barra de aviónica, de modo que todos los equipos se encienden y apagan con un solo interruptor.

\widefigure[0.80\linewidth]{assets/poh/sis-104.png}{Diagrama del sistema eléctrico --- R22 Standard y HP.}
\widefigure[0.80\linewidth]{assets/poh/sis-105.png}{Diagrama del sistema eléctrico --- Alpha, Beta y Beta II anteriores.}
\widefigure[0.80\linewidth]{assets/poh/sis-106.png}{Diagrama del sistema eléctrico --- Beta II posteriores.}
\widefigure[0.78\linewidth]{assets/poh/sis-107.png}{Panel de disyuntores típico (helicópteros anteriores).}
\widefigure[0.78\linewidth]{assets/poh/sis-108.png}{Panel de disyuntores típico (helicópteros posteriores).}

\section{Sistema de iluminación}

Una luz anticolisión roja instalada en el cono de cola se controla con el interruptor \textsf{STROBE}. Las luces de posición, a cada lado de la cabina y en la cola, se controlan con el interruptor \textsf{NAV LTS}. Luces de poste e internas (helicópteros anteriores) o una luz en la parte superior del parabrisas (posteriores) iluminan los instrumentos; la iluminación de instrumentos se activa con el interruptor de luces de navegación y se atenúa con el pomo ubicado sobre él. Una luz de mapa en el techo, montada sobre una rótula, se controla con un interruptor adyacente y puede usarse como iluminación de emergencia del panel.

Dos luces de aterrizaje en la nariz, con ángulos verticales distintos para aumentar el área iluminada, se controlan con un único interruptor en el poste central del cíclico.

\begin{nota}
Las luces de aterrizaje funcionan solo con el interruptor del actuador del embrague en la posición \textsf{ENGAGE}.
\end{nota}

\begin{nota}
Se recomienda operar continuamente las luces de aterrizaje y de posición en vuelo para favorecer la prevención de colisiones.
\end{nota}

\section{Panel de instrumentos}

Los instrumentos primarios estándar son: anemómetro, tacómetro doble de motor y rotor, altímetro, indicador de presión de admisión y brújula magnética. Los instrumentos de motor incluyen amperímetro, presión y temperatura de aceite, temperatura de culata y cantidad de combustible del tanque principal y del auxiliar (si está instalado). También son estándar un reloj, un indicador de temperatura del aire del carburador y un indicador digital de temperatura exterior. El horómetro está en la repisa justo delante del asiento del piloto.

Los pilotos deben familiarizarse con la disposición del panel y el equipamiento de cada helicóptero específico que vuelan; la configuración exacta varía con el equipo opcional y la fecha de fabricación.

\widefigure[0.80\linewidth]{assets/poh/sis-111.png}{Panel de instrumentos típico (helicópteros anteriores). 1~Variómetro; 2~instrumento opcional; 3~anemómetro; 4~tacómetros de motor y rotor; 5~altímetro; 6~instrumento opcional; 7~presión de admisión; 8~luz del actuador del embrague; 9~luz de temperatura de la caja MR; 10~luz de virutas de la caja MR; 11~luz de monóxido de carbono; 12~luz de arranque engranado; 13~luz de virutas de la caja TR; 14~luz de bajo combustible; 15~luz de baja RPM; 16~luz de bajo voltaje ALT; 17~luz de presión de aceite; 18~luz de gobernador desconectado; 19~reloj o temperatura del carburador; 20~temperatura del carburador o reloj; 21~instrumentos de motor; 22~atenuador de luces del panel; 23~luz del freno rotor; 24~interruptor de luces de navegación; 25~interruptor de estrobo; 26~interruptor del actuador del embrague; 27~interruptor del alternador; 28~interruptor de batería; 29~interruptor de encendido; 30~intercomunicador; 31~temperatura exterior/voltímetro; 32~aire de cabina; 33~fricción del cíclico; 34~compensación derecha del cíclico; 35~control de mezcla; 36~interruptor del ELT (opcional); 37~calefacción del carburador.}

\widefigure[0.80\linewidth]{assets/poh/sis-112.png}{Panel de instrumentos opcional (helicópteros anteriores), con horizonte artificial, coordinador de virajes, HSI, radiobaliza y control de esclavización del HSI, además de los elementos del panel estándar.}

\widefigure[0.80\linewidth]{assets/poh/sis-113.png}{Panel de instrumentos típico (helicópteros posteriores). Incluye la luz \textsf{FULL THROTTLE}, la luz del freno rotor, el interruptor maestro de aviónica y la consola lateral del piloto opcional.}

\section{Aviónica opcional}

Existe una amplia gama de aviónica opcional para el R22 y no es práctico describir todos los equipos en este manual. Todos los helicópteros se entregan con los manuales de operación del fabricante de cada equipo instalado, a los que se remite al piloto para las instrucciones detalladas. Para opciones más complejas, como las pantallas primarias de vuelo (PFD), hay información adicional en el suplemento de aviónica opcional (capítulo~\ref{cap:suplementos}). La buena práctica exige familiarizarse con el equipo instalado antes de operar.

\section{Sistema de audio}

Un intercomunicador/sistema de audio activado por voz es estándar y se controla con un pequeño panel sobre el bastidor de aviónica. El pomo \textsf{ICS} controla el volumen del intercomunicador pero no el de la radio. El pomo \textsf{VOX} ajusta el umbral de volumen al que se activa el intercomunicador; girado completamente en sentido horario, se requiere pulsar para activar el intercomunicador. Los sistemas más recientes incluyen una entrada de música en el panel de disyuntores, que se silencia cuando el intercomunicador está activo, al transmitir y durante la recepción de radio.

Las tomas de audífonos están en el techo, cerca de cada asiento. Las empuñaduras del cíclico tienen botones de transmisión e intercomunicador o interruptores tipo gatillo; con el gatillo, el primer retén activa el intercomunicador y el segundo transmite. Hay un botón de intercomunicador adicional en el piso izquierdo o en el soporte del asiento.

Los R22 más antiguos tienen un intercomunicador que opera en modo pulsar para hablar (PTT) o micrófono abierto, seleccionable con un interruptor a la izquierda del poste central del cíclico. En modo PTT se activa con los botones de intercomunicador. Como alternativa al sistema estándar se ofrecen paneles de audio de varios fabricantes; consulte sus instrucciones si hay uno instalado.

\section{Sistema pitot--estático}

El sistema pitot--estático suministra presión de aire al anemómetro, al altímetro y al variómetro. El tubo pitot está en el borde delantero del carenado del mástil. La toma estática está dentro del capó trasero, por dentro de la bisagra de la puerta del capó.

El agua puede drenarse de las líneas retirando los tapones plásticos de drenaje, accesibles por el panel de inspección bajo la cabina. Solo debería ser necesario drenar las líneas si el anemómetro o el altímetro funcionan de forma errática. Las tomas pitot y estática deben inspeccionarse con frecuencia para verificar que no tengan insectos ni otras obstrucciones.

\section{Tacómetro doble}

El tacómetro electrónico doble de motor y rotor es estándar. La señal del tacómetro del motor proviene de los platinos del magneto; la del rotor, de dos emisores magnéticos en la horquilla de accionamiento de la caja principal. Cada tacómetro está en un circuito separado con su propio disyuntor. Con los interruptores de batería y alternador apagados, los tacómetros siguen recibiendo energía de la batería por un circuito de derivación, mientras el interruptor del actuador del embrague esté en \textsf{ENGAGE}.

\begin{nota}
No guarde el helicóptero con el embrague acoplado: los tacómetros quedan energizados y descargarán la batería.
\end{nota}

\section{Luces de advertencia y precaución}

Incluyen: embrague, sobretemperatura de la caja principal, virutas de las cajas principal y de cola, arranque engranado (helicópteros posteriores), bajo combustible, baja RPM, alternador, baja presión de aceite, freno rotor, gobernador desconectado, monóxido de carbono (con calefacción de cabina) y acelerador a fondo (posteriores).

La luz del embrague indica que el actuador está funcionando. La luz y la bocina de baja RPM indican una RPM del rotor de 97\,\% o menos. Las luces de baja presión de aceite y bajo combustible las accionan sensores propios, independientes de los indicadores. La luz del alternador advierte de una posible falla del alternador. La luz de gobernador desconectado indica que el gobernador está apagado.

Los detectores de virutas de las cajas principal y de cola son dispositivos magnéticos ubicados en el tapón de drenaje de cada caja. Cuando las partículas metálicas son atraídas por los imanes, cierran un circuito eléctrico que enciende la luz. Las partículas pueden deberse a un rodamiento o engranaje en falla, por lo que advierten de una falla inminente de la caja. La luz de sobretemperatura de la caja principal la acciona un interruptor térmico cerca del piñón de entrada.

La luz de monóxido de carbono la acciona un sensor sobre la salida de calefacción del piloto. La luz de acelerador a fondo la activa un interruptor en la articulación del acelerador e indica que el motor está cerca del acelerador completamente abierto.

\section{Alertas sonoras}

Todos los R22 tienen una bocina de baja RPM que suena cuando la RPM del rotor está por debajo de 97\,\%; se silencia con el colectivo completamente abajo. En los helicópteros anteriores la produce un parlante en el costado de la consola; en los posteriores, un generador de tono en el sistema de audio la entrega por los audífonos.

Los helicópteros posteriores incluyen una alerta de alta RPM del rotor por los audífonos: un gorjeo (tono alto/bajo) indica que la RPM se acerca a 110\,\%. Un botón de prueba en el panel permite probarla antes o durante el vuelo. Puede haber alertas adicionales según el equipo opcional, como avisos de terreno y tránsito.

\section{Unidad de monitoreo del motor (EMU)}

Los helicópteros posteriores tienen una unidad de monitoreo del motor (\emph{Engine Monitoring Unit}, EMU), un registrador digital dentro de la caja de control del gobernador. La EMU vigila continuamente la RPM del motor y del rotor, la temperatura de aceite, la temperatura de culata, la presión de admisión, la presión ambiente y la temperatura exterior.

El estado de la EMU lo indica una luz ámbar en el compartimiento de equipaje derecho, en el panel trasero del soporte del asiento. La EMU necesita unos diez segundos para completar su autoprueba después de conectar la batería. Luego:

\begin{checklist}
\chk{Luz fija}{Operación normal}
\chk{Destello lento (cada 2 s) o luz apagada}{Falla en los sensores o circuitos de la EMU}
\chk{Destello rápido (4 por segundo)}{La EMU detectó una excedencia}
\end{checklist}

Los datos pueden descargarse a un computador con el software apropiado. Una falla o excedencia debe ser investigada y la indicación reiniciada por un mecánico calificado antes del siguiente vuelo. Las excedencias actuales pueden revisarse (pero no reiniciarse) en un dispositivo móvil Apple con una aplicación disponible en el sitio web de Robinson. La EMU es solo una ayuda de mantenimiento: sigue siendo responsabilidad del piloto informar cualquier excedencia observada.

\section{Calefacción y ventilación de la cabina}

Hay ventilaciones de aire fresco en cada puerta y en la nariz. Las de las puertas se abren y cierran con el pomo junto a la bisagra de la ventilación. En los helicópteros anteriores, empujar el pomo sella y traba la ventilación; en los posteriores, un pomo giratorio la traba y puede usarse para mantenerla parcialmente abierta. Para máxima ventilación, abra las ventilaciones por completo en estacionario, pero solo una pulgada o menos en crucero.

La toma de aire fresco de la nariz se abre tirando del pomo \textsf{CABIN AIR} de la consola. El aire se dirige por la superficie interior del parabrisas para desempañarlo y ventilar.

La calefacción de cabina es opcional: consiste en una camisa de calefacción en el silenciador, una válvula de control en el cortafuego, una rejilla de salida delante de los pedales y los ductos de interconexión. Un pomo en la repisa delante del asiento del piloto acciona la válvula, que dirige el calor hacia la cabina o hacia una descarga al exterior bajo la cabina. Los helicópteros anteriores tienen un soplador eléctrico con interruptor junto al pomo; en los posteriores, el aire se toma del carenado de enfriamiento del motor y no se requiere soplador.

\begin{precaucion}
En caso de fuego de motor en vuelo, el pomo de calefacción debe empujarse a OFF para sellar la cabina respecto del compartimiento del motor.
\end{precaucion}

\begin{nota}
En los helicópteros con soplador de calefacción, haga retirar la camisa de calefacción en las estaciones cálidas para prolongar la vida del silenciador.
\end{nota}

\section{Asientos, cinturones y equipaje}

Los asientos no son ajustables, pero cada helicóptero se entrega con un cojín de respaldo removible que adelanta al piloto. Así los pilotos de menor estatura pueden alcanzar los pedales, la empuñadura del cíclico en su posición más adelantada y los controles de la consola central.

Cada asiento tiene un cinturón combinado de regazo y arnés de hombro con carrete inercial. El carrete normalmente está libre, pero se traba ante un movimiento brusco como el de un accidente. Hay un compartimiento de equipaje bajo cada asiento; los cojines se abaten hacia adelante para acceder.

\section{Tren de aterrizaje}

Se usa un tren de patines. La mayoría de los aterrizajes duros se absorben elásticamente. Sin embargo, en un aterrizaje extremadamente duro, los montantes se abren hacia arriba y hacia afuera a medida que el tubo transversal cede (queda doblado permanentemente) para absorber el impacto. Una leve deformación del tubo transversal es aceptable. Sin embargo, \textbf{una deformación que deje el patín de cola a menos de 34~in del suelo (24~in en el R22 Standard o HP), con el helicóptero vacío sobre pavimento nivelado, exige reemplazar el tubo transversal}.

En la parte inferior de los patines van montadas zapatas de desgaste resistentes a la abrasión. Deben inspeccionarse periódicamente, sobre todo después de aterrizajes en autorrotación. \textbf{Haga reemplazar las zapatas cuando el punto más delgado del área de desgaste mida menos de 0,05~in (1,3~mm).}

\section{Freno rotor}

Si está instalado, el freno rotor va montado en el extremo trasero de la caja principal y se acciona con un cable conectado a una manilla ubicada arriba y detrás del hombro izquierdo del piloto. Para detener el rotor:

\begin{pasos}
\item Después de cortar la mezcla, espere al menos 30 segundos.
\item Tire de la manilla hacia adelante y hacia abajo con fuerza moderada (unas 10~lb).
\item Una vez detenido el rotor, se recomienda usar el freno como freno de estacionamiento enganchando la cadena de bolitas en la ranura del soporte.
\end{pasos}

El freno debe soltarse antes de arrancar el motor. Con el freno aplicado, el motor de arranque queda inhabilitado.

\begin{precaucion}
Aplicar el freno rotor sin esperar al menos 30 segundos después de que se detenga el motor, o con una fuerza que detenga el rotor en menos de 20 segundos, puede dañar las zapatas del freno.
\end{precaucion}

\section{Sistema cebador del motor (opcional)}

El cebador mejora el arranque en frío. La bomba está delante del asiento derecho, cerca del horómetro. Para cebar:

\begin{pasos}
\item Destrabe la manilla de la bomba y bombee lo necesario (normalmente dos o tres carreras). Tire de la manilla lentamente para dar tiempo a que el combustible llene la bomba.
\item Después de cebar, empuje la manilla completamente abajo y trábela.
\end{pasos}

\section{Detector de monóxido de carbono}

El detector de CO se instala si está instalada la calefacción de cabina opcional e indica niveles elevados de CO en la cabina. El CO es un gas inodoro y tóxico presente en los gases de escape, que causa dolor de cabeza, somnolencia y posible pérdida de conciencia. Sus niveles pueden elevarse por una fuga en el escape o por recirculación de los gases durante un estacionario prolongado.

El sistema consiste en un sensor sobre la salida de calefacción del piloto y una luz de precaución. Cada vez que se conecta la energía se hace una autoprueba (la luz destella dos veces). Un destello continuo cada cuatro segundos indica una falla del sensor.

Si la luz se enciende, apague la calefacción y abra las ventilaciones de la nariz y de las puertas. Si está en estacionario, aterrice o pase a vuelo de avance. Si aparecen síntomas de intoxicación (dolor de cabeza, somnolencia, mareo) junto con la luz, aterrice de inmediato. Haga inspeccionar el sistema de escape antes del siguiente vuelo.

Muchos productos químicos pueden dañar el sensor. Evite usar solventes, detergentes o aerosoles cerca de él, y cubra temporalmente con cinta las aberturas superior e inferior de su carcasa al limpiar el interior de la cabina.

\section{Equipo ADS-B}

Los helicópteros posteriores tienen un transpondedor con capacidad ADS-B (\emph{Automatic Dependent Surveillance--Broadcast}), que transmite la posición GPS al control de tránsito aéreo para complementar la información de radar y transpondedor. La capacidad ADS-B ``Out'' es obligatoria en ciertos espacios aéreos. El equipo instalado en fábrica cumple el 14 CFR \S\,91.227. Su operación es casi automática; las fallas se anuncian en el transpondedor o en el GPS.

\begin{nota}
El ADS-B Out instalado en fábrica opera en 1090~MHz, frecuencia aceptada también en la mayoría de los países fuera de los EE.\,UU.
\end{nota}

\begin{nota}
Los controles del transpondedor permiten apagar las transmisiones ADS-B Out, pero como es obligatorio en ciertos espacios aéreos, no debe apagarse salvo instrucción del control de tránsito aéreo.
\end{nota}

El equipo ADS-B también puede recibir información de tránsito de otras aeronaves equipadas y, según el equipo y el país, información adicional de tránsito y meteorología desde estaciones terrestres (ADS-B ``In''), que no es obligatorio. El ADS-B In requiere una pantalla adecuada, como un GPS de mapa móvil o una pantalla multifunción (MFD). El equipo de fábrica recibe ambas frecuencias aprobadas en EE.\,UU. (978 y 1090~MHz).

\section{Transmisor localizador de emergencia (opcional)}

La instalación del ELT consiste en un transmisor con batería interna, una antena externa y un interruptor/anunciador remoto. El transmisor va en el bastidor superior de tubos de acero y es accesible por la puerta superior trasera del capó; el interruptor remoto está a la izquierda del cíclico.

El interruptor del transmisor quedó asegurado en \textsf{AUTO} o \textsf{ARM} al instalarlo y siempre debe estar en esa posición para el vuelo. El interruptor remoto, de tres posiciones con luz indicadora, también debe estar en \textsf{AUTO/ARMED} (posición central). Con ambos así, el ELT empieza a transmitir al sufrir una carga G alta, y la luz roja se enciende mientras transmite.

Mover el interruptor remoto a \textsf{ON} activa el transmisor: úselo si es inminente un aterrizaje de emergencia y hay tiempo. Si el ELT se activa inadvertidamente, use la posición \textsf{RESET} para detener la transmisión y reiniciar la unidad; la luz roja se apagará.

\begin{nota}
Los helicópteros anteriores pueden tener ELT sin interruptor remoto. Para la operación, el mantenimiento y las pruebas requeridas, consulte las instrucciones del fabricante del equipo.
\end{nota}

\section{Soportes de accesorios (opcionales)}

Hay disponible una opción para montar artículos pequeños y portátiles, con barras de montaje delante del asiento del piloto, del copiloto o ambos. Las barras llevan una o más abrazaderas compatibles con diversos soportes comerciales. \textbf{El límite de peso es de 10~lb en total por barra.} Cerca del extremo interior de las barras hay tomas de energía USB y tipo encendedor, protegidas por el disyuntor \textsf{AUX POWER} y fusibles en línea, rotuladas con su voltaje y corriente; también hay abrazaderas para cables y un bolsillo para el cable sobrante.

Las barras, abrazaderas y tomas están aprobadas como parte del diseño tipo, pero los objetos que se fijen son responsabilidad del piloto al mando según las normas de operación aplicables. Asegúrese de que estén firmes y de que no interfieran con los controles ni con el campo visual principal. Pase los cables por las abrazaderas o asegúrelos a la barra con amarras o cinta.

\section{Cámara de cabina (opcional)}

Puede instalarse una cámara de video en el techo de la cabina. Graba video 4K, el audio del intercomunicador y de las comunicaciones y la posición GPS, tanto internamente como en una memoria \emph{flash} removible insertada al frente de la carcasa. La memoria interna conserva solo el video reciente y no es accesible al usuario. La grabación empieza automáticamente al conectar la batería y se detiene al desconectarla.

La grabación en la memoria puede detenerse o silenciarse el audio con los interruptores de grabación y audio al frente de la carcasa. No retire la memoria mientras graba, pues se corrompe el archivo; con la batería conectada, detenga primero la grabación. Una luz azul destellante indica grabación en la memoria; una luz verde fija, cámara energizada y funcionando; una luz ámbar destellante, una falla interna (puede no estar grabándose).

El video se graba en archivos secuenciales de 4~GB (unos 25 minutos cada uno) llamados \texttt{HELICAM\_xxxx.MP4}; la posición y altitud GPS pueden mostrarse en el video y se guardan aparte en archivos \texttt{HELICAM\_xxxx.GPX}. Una memoria de 128~GB (la suministrada) graba unas 10 horas; cuando se llena, el archivo más antiguo se sobrescribe.

\begin{nota}
Las memorias deben cumplir los criterios de la guía del usuario de la cámara (\emph{Cockpit Camera User Guide}, disponible en el sitio web de Robinson) para funcionar de forma confiable.
\end{nota}
